\section{Contexte et objectifs}

% ── Contexte / problématique ────────────────────────────────────────────────
\begin{frame}{Contexte et problématique}
  \textbf{Classer des candidatures académiques} mobilise plusieurs critères et plusieurs
  décideurs : un processus difficile à rendre \alert{objectif, cohérent et transparent}.

  \vspace{0.6em}
  \begin{itemize}
    \item Subjectivité des jugements, pondération implicite des critères.
    \item Les critères pertinents \alert{varient d'une formation à l'autre}
          (parcours, compétences, résultats, expériences).
    \item \textbf{Besoin} : un outil \alert{paramétrable}, transparent et qui s'améliore
          avec les décisions passées.
  \end{itemize}

  \vspace{0.6em}
  \begin{block}{Question centrale}
    Comment combiner la rigueur d'une méthode multicritère et la capacité prédictive de
    l'apprentissage automatique, tout en restant \emph{générique}, \emph{explicable} et
    \emph{rigoureusement évalué} ?
  \end{block}
\end{frame}

% ── Objectifs ───────────────────────────────────────────────────────────────
\begin{frame}{Objectifs}
  Concevoir un système hybride d'aide à la décision qui :
  \begin{enumerate}
    \item pondère les critères par \textbf{AHP} \footnotesize(\emph{Analytic Hierarchy
          Process})\normalsize{} avec contrôle de cohérence ;
    \item classe les candidats par \textbf{TOPSIS} \footnotesize(\emph{Technique for Order
          of Preference by Similarity to Ideal Solution})\normalsize{} ;
    \item prédit une classe via un \textbf{modèle d'apprentissage} appris sur une vérité
          terrain, évalué \alert{sans fuite de données} ;
    \item combine les deux par une \textbf{fusion hybride} pondérée ;
    \item \textbf{justifie} chaque décision : contributions par critère (TOPSIS) et par
          variable via \textbf{SHAP} \footnotesize(\emph{SHapley Additive exPlanations})\normalsize.
  \end{enumerate}
\end{frame}

% ── État de l'art ───────────────────────────────────────────────────────────
\section{État de l'art}

% ── Diagramme PRISMA ────────────────────────────────────────────────────────
\begin{frame}{Méthodologie de la revue (PRISMA)}
  \begin{center}
  \begin{tikzpicture}[
    box/.style={rectangle, draw=primary, fill=softgray, rounded corners,
                text width=5.4cm, align=center, inner sep=3pt, font=\scriptsize, minimum height=0.7cm},
    excl/.style={rectangle, draw=gray, fill=white, rounded corners,
                text width=3.8cm, align=center, inner sep=3pt, font=\scriptsize, minimum height=0.7cm},
    arr/.style={-{Latex[length=1.6mm]}, thick, draw=primary},
    earr/.style={-{Latex[length=1.6mm]}, semithick, draw=gray},
    node distance=0.45cm and 1.2cm
  ]
    \node[box] (id) {Notices identifiées (5 bases) : \textbf{n = 142}};
    \node[box, below=of id] (dedup) {Après suppression des doublons : \textbf{n = 119}};
    \node[box, below=of dedup] (full) {Évalués en texte intégral : \textbf{n = 64}};
    \node[box, below=of full, fill=primary, text=white] (incl) {Inclus dans la synthèse : \textbf{n = 56} (dont 12 principaux)};

    \node[excl, right=of dedup] (ex1) {Doublons : \textbf{23}};
    \node[excl, right=of full] (ex2) {Exclus (titre/résumé, non récupérés, hors champ) : \textbf{63}};

    \draw[arr] (id)--(dedup);
    \draw[arr] (dedup)--(full);
    \draw[arr] (full)--(incl);
    \draw[earr] (dedup.east)--(ex1.west);
    \draw[earr] (full.east)--(ex2.west);
  \end{tikzpicture}
  \end{center}
  \footnotesize Sources : IEEE Xplore, ScienceDirect, SpringerLink, ACM DL, Google Scholar.
\end{frame}

% ── Synthèse / positionnement ───────────────────────────────────────────────
\begin{frame}{Positionnement}
  \vspace{0.3em}
  \begin{columns}[T]
    \begin{column}{0.5\textwidth}
      \textbf{Ce que dit la littérature}
      \begin{itemize}\small
        \item AHP--TOPSIS : mûr pour la sélection, mais \alert{normatif} (n'exploite pas
              l'historique).
        \item Hybridation \textbf{MCDM} \scriptsize(aide à la décision multicritère)\small{}
              + ML \scriptsize(apprentissage automatique)\small{} : prometteuse et générique.
        \item Explicabilité (SHAP) : exigence croissante.
      \end{itemize}
    \end{column}
    \begin{column}{0.5\textwidth}
      \textbf{Verrou peu traité}
      \begin{itemize}\small
        \item La \alert{validité de l'évaluation} du modèle (risque de \emph{fuite de
              données}) est rarement discutée dans les travaux appliqués.
      \end{itemize}
      \vspace{0.4em}
      \footnotesize Inspirations principales : Poongothai \textit{et al.} (2025),
      James \textit{et al.} (2023).
    \end{column}
  \end{columns}
\end{frame}
